\subsection{Morphisms of locally constant sheaves on a locally connected space}

Let B be a topological space, let \((\mathrm{E},p)\) and \((\mathrm{E}^{\prime},p^{\prime})\) be B-spaces. Denote by \(\mathcal{M} = \underline{\mathcal{M}or_{\mathrm{B}}(\mathrm{E};\mathrm{E}^{\prime})}\) the sheaf on B of B-morphisms from \((\mathrm{E},p)\) to \((\mathrm{E}^{\prime},p^{\prime})\) (I, p. 45, example 4). If U is an open set of B and b a point of U, we denote by \(\theta_{b,\mathrm{U}}\colon\mathcal{M}(\mathrm{U})\to\mathcal{C}(\mathrm{E}_{b};\mathrm{E}_{b}^{\prime})\) the canonical map obtained by passing to fibers at b. We denote by \(\theta_{b}\colon\mathcal{M}_{b}\to\mathcal{C}(\mathrm{E}_{b};\mathrm{E}_{b}^{\prime})\) the unique map such that \(\theta_{b,\mathrm{U}}\) is the composite of \(\theta_{b}\) and the canonical map \(\mathcal{M}(\mathrm{U})\to\mathcal{M}_{b}\) for every open set U of B containing b (E, III, p. 62).

Let also \(\mathcal{I} = \underline{\mathcal{I}som}_{\mathrm{B}}(\mathrm{E};\mathrm{E}^{\prime})\) be the sheaf on B of B-isomorphisms from \((\mathrm{E},p)\) to \((\mathrm{E}^{\prime},p^{\prime})\). Denote by \(i\colon \mathcal{I}\to \mathcal{M}\) the canonical morphism. For every \(b\in \mathrm{B}\), the map \(\theta_{b}\circ i_{b}\) induces a map \(\theta_b^\prime\) of \(\mathcal{I}_b\) into the set of continuous bijections of \(\mathrm{E}_b\) onto \(\mathrm{E}_b^\prime\).

PROPOSITION 12. --- Suppose that the space B is locally connected and that the B-spaces E and E' are coverings. The sheaf \(\mathcal{M}or_{\mathrm{B}}(\mathrm{E};\mathrm{E}')\) is then locally constant and, for every \(b\in\mathrm{B}\), the map \(\theta_{b}\) is a bijection from its stalk at \(b\) onto the set of maps from \(\mathrm{E}_{b}\) to \(\mathrm{E}_{b}^{\prime}\).

Similarly, the sheaf \(\mathcal{I}\text{som}_{\mathrm{B}}(\mathrm{E};\mathrm{E}^{\prime})\) is locally constant and, for every point \(b\) of B, the map \(\theta_{b}^{\prime}\) is a bijection from its stalk at \(b\) onto the set of bijections of \(\mathrm{E}_{b}\) onto \(\mathrm{E}_{b}^{\prime}\).

We shall denote by \(\mathcal{M} = \underline{\mathcal{M}or}_{\mathrm{B}}(\mathrm{E};\mathrm{E}^{\prime})\) and \(\mathcal{I} = \underline{\mathcal{Isom}}_{\mathrm{B}}(\mathrm{E};\mathrm{E}^{\prime})\). The assertions to be proved are local in nature on B; we may therefore assume that \(\mathrm{E} = \mathrm{B}\times \mathrm{F}\) and \(\mathrm{E}' = \mathrm{B}\times \mathrm{F}'\) are trivial coverings, where F and \(\mathrm{F}'\) are discrete topological spaces.

First, let us prove that, for every connected open subset U of B and every point b of U, the map \(\theta_{b,U}\) is a bijection. Let \(f\colon U\times F\to U\times F'\) be a morphism of U-spaces; then, \(\theta_{b,U}(f)\) is the map \(y\mapsto\mathrm{pr}_{2}(f(b,y))\) from F to \(F'\). For all \(y\in F\), the map \(x\mapsto\mathrm{pr}_{2}(f(x,y))\) is a continuous map from the connected space U into the discrete space \(F'\), hence is constant, equal to \(\mathrm{pr}_{2}(f(b,y))\). One thus has \(f(x,y)=(x,\theta_{b,U}(f)(y))\). It follows that \(\theta_{b,U}\) is a bijection;

its inverse bijection associates to every map \(\varphi\colon\mathrm{F}\to\mathrm{F}^{\prime}\) the U-morphism from \(\mathrm{U}\times\mathrm{F}\) to \(\mathrm{U}\times\mathrm{F}^{\prime}\) given by \((x,y)\mapsto(x,\varphi(y))\).

By hypothesis, every point \(b \in B\) admits a base of connected open neighborhoods; the map \(\theta_{b}\) is therefore a bijection. The maps \(\theta_{b,\mathrm{U}}^{-1}\), for U a nonempty connected open subset of B and b an arbitrary point of B, define a morphism of presheaves (with respect to the base of connected open subsets of B) from the constant sheaf with stalk-type \(F^{\prime F}\) to the sheaf \(\mathcal{M}\). By what precedes, this morphism induces a bijection on stalks; it is therefore an isomorphism.

The assertions concerning the sheaf \(\mathcal{I}\) are proved similarly.

COROLLARY 1. --- Let B be a topological space and A a subspace of B. Suppose that the spaces A and B are locally connected. Let E and E' be coverings of B. Then the canonical morphisms \(\psi\colon\mathcal{M}or_{B}(E;E')_{A}\to\mathcal{M}or_{A}(E_{A};E'_{A})\) and \(\psi'\colon\mathcal{I}som_{B}(E;E')_{A}\to\mathcal{I}som_{A}(E_{A};E'_{A})\) (I, p. 45, example 4) are isomorphisms.

For every point \(a \in \mathrm{A}\), it follows from the preceding proposition and from example 2 of I, p. 63, applied to the spaces \(\{a\}\), A, B and to the canonical injections, that the canonical morphism \(\psi\) induces, by passage to stalks, the identity on \(\mathrm{E}_a^{\mathrm{E}_a'}\). It is therefore an isomorphism. The fact that \(\psi'\) is an isomorphism is proved in an analogous manner.

COROLLARY 2. --- Let B be a topological space and A a subspace of B. Suppose that the spaces A and B are locally connected and that the pair (B, A) has property (PCV) of I, p. 37. Let E and E' be coverings of B, and denote their projections by p and p'. Let \(g: \overline{p}^{1}(A) \to (\overline{p}')^{1}(A)\) be an A-morphism (resp. an A-isomorphism). There exists a neighborhood U of A in B and a U-morphism (resp. a U-isomorphism) \(f: \overline{p}^{1}(U) \to (\overline{p}')^{1}(U)\) such that \(f_{A} = g\).

Let us keep the notation of Corollary 1. By that same corollary, an A-morphism \(g \colon \mathrm{E}_{\mathrm{A}} \to \mathrm{E}_{\mathrm{A}}^{\prime}\) is identified with a section \(s_0\) over A of the étale space \(\mathrm{E}_{\mathcal{M}}\) associated with \(\mathcal{M}\). By the hypothesis made on the pair (B, A) and Lemma 3 of I, p. 39, there exists an open neighborhood U of A in B and a continuous section \(s\) of \(\mathrm{E}_{\mathcal{M}}\) above U extending \(s_0\). This section \(s\) is identified with a U-morphism \(f \colon \mathrm{E}_{\mathrm{U}} \to \mathrm{E}_{\mathrm{U}}^{\prime}\) extending \(g\).

The case where g is an A-isomorphism is handled analogously by considering, in place of the sheaf M, the sheaf S.

